% Vendored verbatim excerpts from The Stacks Project, master branch,
% retrieved 2026-06-04. Source URLs and exact line ranges noted below.
% Do NOT edit content — only headers/comments added by the retriever.
%
% =====================================================================
% EXCERPT 1 — Algebra chapter, "Henselian local rings" section
% Source: https://raw.githubusercontent.com/stacks/stacks-project/master/algebra.tex
% Verbatim copy of lines 42066–42433.
% Contains: section header (label section-henselian),
%           definition-henselian             [tag 04GF],
%           lemma-uniqueness                 [in proof of 04GG (1)=>(8)],
%           lemma-characterize-henselian     [tag 04GG] + full proof,
%           lemma-finite-over-henselian      [tag 04GH] + full proof,
%           lemma-mop-up                     [tag 04GJ] + proof.
% The section label "section-henselian" itself is tag 04GE.
% =====================================================================

\section{Henselian local rings}
\label{section-henselian}

\noindent
In this section we discuss a bit the notion of a henselian local ring.
Let $(R, \mathfrak m, \kappa)$ be a local ring.
For $a \in R$ we denote $\overline{a}$ the image of $a$ in $\kappa$.
For a polynomial $f \in R[T]$ we often denote $\overline{f}$
the image of $f$ in $\kappa[T]$.
Given a polynomial $f \in R[T]$ we denote $f'$ the derivative
of $f$ with respect to $T$. Note that $\overline{f}' = \overline{f'}$.

\begin{definition}
\label{definition-henselian}
Let $(R, \mathfrak m, \kappa)$ be a local ring.
\begin{enumerate}
\item We say $R$ is {\it henselian} if for every monic $f \in R[T]$ and
every root $a_0 \in \kappa$ of $\overline{f}$ such that
$\overline{f'}(a_0) \not = 0$
there exists an $a \in R$ such that $f(a) = 0$ and
$a_0 = \overline{a}$.
\item We say $R$ is {\it strictly henselian} if $R$ is henselian
and its residue field is separably algebraically closed.
\end{enumerate}
\end{definition}

\noindent
Note that the condition $\overline{f'}(a_0) \not = 0$ is equivalent to the
condition that $a_0$ is a simple root of the polynomial $\overline{f}$.
In fact, it implies that the lift $a \in R$, if it exists, is unique.

\begin{lemma}
\label{lemma-uniqueness}
Let $(R, \mathfrak m, \kappa)$ be a local ring.
Let $f \in R[T]$. Let $a, b \in R$ such that $f(a) = f(b) = 0$,
$a = b \bmod \mathfrak m$, and $f'(a) \not \in \mathfrak m$.
Then $a = b$.
\end{lemma}

\begin{proof}
Write $f(x + y) - f(x) = f'(x)y + g(x, y) y^2$ in $R[x, y]$ (this is possible
as one sees by expanding $f(x + y)$; details omitted).
Then we see that $0 = f(b) - f(a) = f(a + (b - a)) - f(a) =
f'(a)(b - a) + c (b - a)^2$ for some $c \in R$. By assumption
$f'(a)$ is a unit in $R$. Hence $(b - a)(1 + f'(a)^{-1}c(b - a)) = 0$.
By assumption $b - a \in \mathfrak m$, hence $1 + f'(a)^{-1}c(b - a)$
is a unit in $R$. Hence $b - a = 0$ in $R$.
\end{proof}

\noindent
Here is the characterization of henselian local rings.

\begin{lemma}
\label{lemma-characterize-henselian}
\begin{slogan}
Characterizations of henselian local rings
\end{slogan}
Let $(R, \mathfrak m, \kappa)$ be a local ring.
The following are equivalent
\begin{enumerate}
\item $R$ is henselian,
\item for every $f \in R[T]$ and every root $a_0 \in \kappa$
of $\overline{f}$ such that $\overline{f'}(a_0) \not = 0$
there exists an $a \in R$ such that $f(a) = 0$ and
$a_0 = \overline{a}$,
\item for any monic $f \in R[T]$ and any factorization
$\overline{f} = g_0 h_0$ with $\gcd(g_0, h_0) = 1$ there
exists a factorization $f = gh$ in $R[T]$ such that
$g_0 = \overline{g}$ and $h_0 = \overline{h}$,
\item for any monic $f \in R[T]$ and any factorization
$\overline{f} = g_0 h_0$ with $\gcd(g_0, h_0) = 1$ there
exists a factorization $f = gh$ in $R[T]$ such that
$g_0 = \overline{g}$ and $h_0 = \overline{h}$ and moreover
$\deg_T(g) = \deg_T(g_0)$,
\item for any $f \in R[T]$ and any factorization
$\overline{f} = g_0 h_0$ with $\gcd(g_0, h_0) = 1$ there
exists a factorization $f = gh$ in $R[T]$ such that
$g_0 = \overline{g}$ and $h_0 = \overline{h}$,
\item for any $f \in R[T]$ and any factorization
$\overline{f} = g_0 h_0$ with $\gcd(g_0, h_0) = 1$ there
exists a factorization $f = gh$ in $R[T]$ such that
$g_0 = \overline{g}$ and $h_0 = \overline{h}$ and
moreover $\deg_T(g) = \deg_T(g_0)$,
\item for any \'etale ring map $R \to S$ and prime $\mathfrak q$ of $S$
lying over $\mathfrak m$ with $\kappa = \kappa(\mathfrak q)$
there exists a retraction $\tau : S \to R$ of $R \to S$,
\item for any \'etale ring map $R \to S$ and prime $\mathfrak q$ of $S$
lying over $\mathfrak m$ with $\kappa = \kappa(\mathfrak q)$
there exists a unique retraction $\tau : S \to R$ of $R \to S$ such that
$\mathfrak q = \tau^{-1}(\mathfrak m)$,
\item any finite $R$-algebra is a product of local rings,
\item any finite $R$-algebra is a finite product of local rings,
\item any finite type $R$-algebra $S$ can be written as
$A \times B$ with $R \to A$ finite
and $R \to B$ not quasi-finite at any prime lying over $\mathfrak m$,
\item any finite type $R$-algebra $S$ can be written as
$A \times B$ with $R \to A$ finite
such that each irreducible component of $\Spec(B \otimes_R \kappa)$
has dimension $\geq 1$, and
\item any quasi-finite $R$-algebra $S$ can be written as
$S = A \times B$ with $R \to A$ finite such that $B \otimes_R \kappa = 0$.
\end{enumerate}
\end{lemma}

\begin{proof}
Here is a list of the easier implications:
\begin{enumerate}
\item 2$\Rightarrow$1 because in (2) we consider all polynomials and
in (1) only monic ones,
\item 5$\Rightarrow$3 because in (5) we consider all polynomials and
in (3) only monic ones,
\item 6$\Rightarrow$4 because in (6) we consider all polynomials and
in (4) only monic ones,
\item 4$\Rightarrow$3 is obvious,
\item 6$\Rightarrow$5 is obvious,
\item 8$\Rightarrow$7 is obvious,
\item 10$\Rightarrow$9 is obvious,
\item 11$\Leftrightarrow$12 by definition of being quasi-finite at a prime,
\item 11$\Rightarrow$13 by definition of being quasi-finite,
\end{enumerate}

\noindent
Proof of 1$\Rightarrow$8. Assume (1).
Let $R \to S$ be \'etale, and let $\mathfrak q \subset S$
be a prime ideal such that $\kappa(\mathfrak q) \cong \kappa$. By
Proposition \ref{proposition-etale-locally-standard}
we can find a $g \in S$, $g \not \in \mathfrak q$ such that
$R \to S_g$ is standard \'etale. After replacing $S$ by $S_g$ we may assume
that $S = R[t]_g/(f)$ is standard \'etale (details omitted).
Since the prime $\mathfrak q$ has residue field $\kappa$ it
corresponds to a root $a_0$ of $\overline{f}$ which is not a
root of $\overline{g}$. By definition of a standard \'etale algebra
this also means that $\overline{f'}(a_0) \not = 0$.
Since also $f$ is monic by definition of a standard \'etale algebra again we
may use that $R$ is henselian to conclude that there exists an $a \in R$
with $a_0 = \overline{a}$ such that $f(a) = 0$. This implies that
$g(a)$ is a unit of $R$ and we obtain the desired map
$\tau : S = R[t]_g/(f) \to R$ by the rule $t \mapsto a$. By construction
$\tau^{-1}(\mathfrak m) = \mathfrak q$. By Lemma \ref{lemma-uniqueness}
the map $\tau$ is unique. This proves (8) holds.

\medskip\noindent
Proof of 7$\Rightarrow$8. (This is really unimportant and should be
skipped.) Assume (7) holds and assume $R \to S$ is \'etale.
Let $\mathfrak q_1, \ldots, \mathfrak q_r$ be
the other primes of $S$ lying over $\mathfrak m$.
Then we can find a $g \in S$, $g \not \in \mathfrak q$ and
$g \in \mathfrak q_i$ for $i = 1, \ldots, r$.
Namely, we can argue that
$\bigcap_{i=1}^{r} \mathfrak{q}_{i} \not\subset \mathfrak{q}$
since otherwise
$\mathfrak{q}_{i} \subset \mathfrak{q}$
for some $i$, but this cannot happen as the fiber of an
\'etale morphism is discrete (use Lemma \ref{lemma-etale-over-field}
for example).
Apply (7) to the \'etale ring map
$R \to S_g$ and the prime $\mathfrak qS_g$. This gives a retraction
$\tau_g : S_g \to R$ such that the composition $\tau : S \to S_g \to R$
has the property $\tau^{-1}(\mathfrak m) = \mathfrak q$.
Details omitted.

\medskip\noindent
Proof of 8$\Rightarrow$11. Assume (8) and let $R \to S$ be a finite type
ring map. Apply
Lemma \ref{lemma-etale-makes-quasi-finite-finite}.
We find an \'etale ring map $R \to R'$ and a prime $\mathfrak m' \subset R'$
lying over $\mathfrak m$ with $\kappa = \kappa(\mathfrak m')$
such that $R' \otimes_R S = A' \times B'$ with $A'$ finite over $R'$
and $B'$ not quasi-finite over $R'$ at any prime lying over $\mathfrak m'$.
Apply (8) to get a retraction $\tau : R' \to R$ with
$\mathfrak m = \tau^{-1}(\mathfrak m')$. Then use that
$$
S = (S \otimes_R R') \otimes_{R', \tau} R
= (A' \times B') \otimes_{R', \tau} R
= (A' \otimes_{R', \tau} R)  \times  (B' \otimes_{R', \tau} R)
$$
which gives a decomposition as in (11).

\medskip\noindent
Proof of 8$\Rightarrow$10. Assume (8) and let $R \to S$ be a finite
ring map. Apply
Lemma \ref{lemma-etale-makes-quasi-finite-finite}.
We find an \'etale ring map $R \to R'$ and a prime $\mathfrak m' \subset R'$
lying over $\mathfrak m$ with $\kappa = \kappa(\mathfrak m')$
such that $R' \otimes_R S = A'_1 \times \ldots \times A'_n \times B'$
with $A'_i$ finite over $R'$ having exactly one prime over $\mathfrak m'$
and $B'$ not quasi-finite over $R'$ at any prime lying over $\mathfrak m'$.
Apply (8) to get a retraction $\tau : R' \to R$ with
$\mathfrak m' = \tau^{-1}(\mathfrak m)$. Then we obtain
\begin{align*}
S & = (S \otimes_R R') \otimes_{R', \tau} R \\
& = (A'_1 \times \ldots \times A'_n \times B') \otimes_{R', \tau} R \\
& = (A'_1 \otimes_{R', \tau} R)  \times
\ldots \times (A'_1 \otimes_{R', \tau} R) \times
(B' \otimes_{R', \tau} R) \\
& = A_1 \times \ldots \times A_n \times B
\end{align*}
The factor $B$ is finite over $R$ but $R \to B$
is not quasi-finite at any prime lying over $\mathfrak m$. Hence
$B = 0$. The factors $A_i$ are finite $R$-algebras having exactly
one prime lying over $\mathfrak m$, hence they are local rings.
This proves that $S$ is a finite product of local rings.

\medskip\noindent
Proof of 9$\Rightarrow$10. This holds because if $S$ is finite over the local
ring $R$, then it has at most finitely many maximal ideals. Namely, by
going up for $R \to S$ the maximal ideals of $S$ all lie over $\mathfrak m$,
and $S/\mathfrak mS$ is Artinian hence has finitely many primes.

\medskip\noindent
Proof of 10$\Rightarrow$1. Assume (10). Let $f \in R[T]$ be a monic
polynomial and $a_0 \in \kappa$ a simple root of $\overline{f}$.
Then $S = R[T]/(f)$ is a finite $R$-algebra. Applying (10)
we get $S = A_1 \times \ldots \times A_r$ is a finite product of
local $R$-algebras. In particular we see that
$S/\mathfrak mS = \prod A_i/\mathfrak mA_i$ is the decomposition
of $\kappa[T]/(\overline{f})$ as a product of local rings.
This means that one of the factors, say $A_1/\mathfrak mA_1$
is the quotient $\kappa[T]/(\overline{f}) \to \kappa[T]/(T - a_0)$.
Since $A_1$ is a summand of the finite free $R$-module $S$ it
is a finite free $R$-module itself. As $A_1/\mathfrak mA_1$ is a
$\kappa$-vector space of dimension 1 we see that $A_1 \cong R$ as an
$R$-module. Clearly this means that $R \to A_1$ is an isomorphism.
Let $a \in R$ be the image of $T$ under the map
$R[T] \to S \to A_1 \to R$. Then $f(a) = 0$ and $\overline{a} = a_0$
as desired.

\medskip\noindent
Proof of 13$\Rightarrow$1. Assume (13). Let $f \in R[T]$ be a monic
polynomial and $a_0 \in \kappa$ a simple root of $\overline{f}$.
Then $S_1 = R[T]/(f)$ is a finite $R$-algebra. Let $g \in R[T]$
be any element such that $\overline{g} = \overline{f}/(T - a_0)$.
Then $S = (S_1)_g$ is a quasi-finite $R$-algebra such that
$S \otimes_R \kappa \cong \kappa[T]_{\overline{g}}/(\overline{f})
\cong \kappa[T]/(T - a_0) \cong \kappa$.
Applying (13) to $S$ we get $S = A \times B$ with $A$ finite over $R$ and
$B \otimes_R \kappa = 0$. In particular we see that
$\kappa \cong S/\mathfrak mS = A/\mathfrak mA$.
Since $A$ is a summand of the flat $R$-algebra $S$ we see
that it is finite flat, hence free over $R$.
As $A/\mathfrak mA$ is a
$\kappa$-vector space of dimension 1 we see that $A \cong R$ as an
$R$-module. Clearly this means that $R \to A$ is an isomorphism.
Let $a \in R$ be the image of $T$ under the map
$R[T] \to S \to A \to R$. Then $f(a) = 0$ and $\overline{a} = a_0$
as desired.

\medskip\noindent
Proof of 8$\Rightarrow$2. Assume (8). Let $f \in R[T]$ be any
polynomial and let $a_0 \in \kappa$ be a simple root. Then
the algebra $S = R[T]_{f'}/(f)$ is \'etale over $R$.
Let $\mathfrak q \subset S$ be the prime
generated by $\mathfrak m$ and $T - b$ where $b \in R$ is any
element such that $\overline{b} = a_0$. Apply (8) to $S$ and $\mathfrak q$
to get $\tau : S \to R$.
Then the image $\tau(T) = a \in R$ works in (2).

\medskip\noindent
At this point we see that (1), (2), (7), (8), (9), (10), (11), (12), (13) are
all equivalent. The weakest assertion of (3), (4), (5) and (6)
is (3) and the strongest is (6). Hence we still have to prove that
(3) implies (1) and (1) implies (6).

\medskip\noindent
Proof of 3$\Rightarrow$1. Assume (3). Let $f \in R[T]$ be monic and
let $a_0 \in \kappa$ be a simple root of $\overline{f}$. This gives
a factorization $\overline{f} = (T - a_0)h_0$ with $h_0(a_0) \not = 0$,
so $\gcd(T - a_0, h_0) = 1$. Apply (3) to get a factorization
$f = gh$ with $\overline{g} = T - a_0$ and $\overline{h} = h_0$.
Set $S = R[T]/(f)$ which is a finite free $R$-algebra. We will write
$g$, $h$ also for the images of $g$ and $h$ in $S$. Then
$gS + hS = S$ by
Nakayama's Lemma \ref{lemma-NAK}
as the equality holds modulo $\mathfrak m$. Since $gh = f = 0$ in $S$
this also implies that $gS \cap hS = 0$. Hence by the Chinese Remainder
theorem we obtain $S = S/(g) \times S/(h)$. This implies that
$A = S/(g)$ is a summand of a finite free $R$-module, hence finite
free. Moreover, the rank of $A$ is $1$ as
$A/\mathfrak mA = \kappa[T]/(T - a_0)$. Thus the map $R \to A$
is an isomorphism. Setting $a \in R$ equal to the image of $T$
under the maps $R[T] \to S \to A \to R$ gives an element of $R$
with $f(a) = 0$ and $\overline{a} = a_0$.

\medskip\noindent
Proof of 1$\Rightarrow$6. Assume (1) or equivalently all of
(1), (2), (7), (8), (9), (10), (11), (12), (13).
Let $f \in R[T]$ be a polynomial.
Suppose that $\overline{f} = g_0h_0$ is a factorization with
$\gcd(g_0, h_0) = 1$. We may and do assume that $g_0$ is monic.
Consider $S = R[T]/(f)$. Because we
have the factorization we see that the coefficients of
$f$ generate the unit ideal in $R$.
This implies that $S$ has finite fibres over $R$, hence is
quasi-finite over $R$. It also implies that $S$ is flat over $R$ by
Lemma \ref{lemma-grothendieck-general}.
Combining (13) and (10) we may write
$S = A_1 \times \ldots \times A_n \times B$
where each $A_i$ is local and finite over $R$, and
$B \otimes_R \kappa = 0$. After reordering the factors $A_1, \ldots, A_n$
we may assume that
$$
\kappa[T]/(g_0) =
A_1/\mathfrak m A_1 \times \ldots \times A_r/\mathfrak mA_r,
\ \kappa[T]/(h_0) =
A_{r + 1}/\mathfrak mA_{r + 1} \times \ldots \times A_n/\mathfrak mA_n
$$
as quotients of $\kappa[T]$. The finite flat $R$-algebra
$A = A_1 \times \ldots \times A_r$ is free as an $R$-module, see
Lemma \ref{lemma-finite-flat-local}.
Its rank is $\deg_T(g_0)$. Let $g \in R[T]$ be the characteristic polynomial
of the $R$-linear operator $T : A \to A$. Then $g$ is a monic polynomial
of degree $\deg_T(g) = \deg_T(g_0)$ and moreover $\overline{g} = g_0$.
By Cayley-Hamilton
(Lemma \ref{lemma-charpoly})
we see that $g(T_A) = 0$ where $T_A$ indicates
the image of $T$ in $A$. Hence we obtain a well defined surjective map
$R[T]/(g) \to A$ which is an isomorphism by
Nakayama's Lemma \ref{lemma-NAK}. The map $R[T] \to A$ factors
through $R[T]/(f)$ by construction hence we may write $f = gh$ for
some $h$. This finishes the proof.
\end{proof}

\begin{lemma}
\label{lemma-finite-over-henselian}
Let $(R, \mathfrak m, \kappa)$ be a henselian local ring.
\begin{enumerate}
\item If $R \to S$ is a finite ring map then $S$ is
a finite product of henselian local rings each finite over $R$.
\item If $R \to S$ is a finite ring map and $S$ is local, then
$S$ is a henselian local ring and $R \to S$ is a (finite) local ring map.
\item If $R \to S$ is a finite type ring map, and $\mathfrak q$ is
a prime of $S$ lying over $\mathfrak m$ at which $R \to S$ is quasi-finite,
then $S_{\mathfrak q}$ is henselian and finite over $R$.
\item If $R \to S$ is quasi-finite then $S_{\mathfrak q}$ is henselian
and finite over $R$ for every prime $\mathfrak q$ lying over $\mathfrak m$.
\end{enumerate}
\end{lemma}

\begin{proof}
Part (2) implies part (1) since $S$ as in part (1) is a finite product
of its localizations at the primes lying over $\mathfrak m$ by
Lemma \ref{lemma-characterize-henselian} part (10).
Part (2) also follows from Lemma \ref{lemma-characterize-henselian} part (10)
since any finite $S$-algebra is also a finite $R$-algebra (of course
any finite ring map between local rings is local).

\medskip\noindent
Let $R \to S$ and $\mathfrak q$ be as in (3). Write $S = A \times B$
with $A$ finite over $R$ and $B$ not quasi-finite over $R$ at any
prime lying over $\mathfrak m$, see
Lemma \ref{lemma-characterize-henselian} part (11).
Hence $S_\mathfrak q$ is a localization of $A$ at a maximal ideal
and we deduce (3) from (1). Part (4) follows from part (3).
\end{proof}

\begin{lemma}
\label{lemma-mop-up}
Let $(R, \mathfrak m, \kappa)$ be a henselian local ring.
Any finite type $R$-algebra $S$ can be written as
$S = A_1 \times \ldots \times A_n \times B$ with $A_i$ local
and finite over $R$ and $R \to B$ not quasi-finite at any
prime of $B$ lying over $\mathfrak m$.
\end{lemma}

\begin{proof}
This is a combination of parts (11) and (10) of
Lemma \ref{lemma-characterize-henselian}.
\end{proof}

% =====================================================================
% EXCERPT 2 — Local Cohomology chapter, "Cohomological dimension" section
% Source: https://raw.githubusercontent.com/stacks/stacks-project/master/local-cohomology.tex
% Verbatim copy of lines 527–539.
% Contains: lemma-cd-local [tag 0DXB] + full proof.
%
% !!! TAG MISMATCH WARNING !!!
% The directive describes tag 0DXB as "residue-product Hensel-lift
% reduction (orthogonal idempotents lifting from S/mS to S)". This is
% NOT what tag 0DXB actually says. Tag 0DXB is a lemma about
% cohomological dimension localising at primes (cd(A,I) = max
% cd(A_p, I_p)). It lives in chapter "Local Cohomology", section
% "Cohomological dimension", and has no connection to henselian rings.
% See the pointer file `stacks-04gg.md` (section "Caveats") for the
% likely intended tags.
% =====================================================================

\begin{lemma}
\label{lemma-cd-local}
Let $I \subset A$ be a finitely generated ideal of a ring $A$.
Then $\text{cd}(A, I) = \max \text{cd}(A_\mathfrak p, I_\mathfrak p)$.
\end{lemma}

\begin{proof}
Let $Y = V(I)$ and $Y' = V(I_\mathfrak p) \subset \Spec(A_\mathfrak p)$.
Recall that
$R\Gamma_Y(A) \otimes_A A_\mathfrak p = R\Gamma_{Y'}(A_\mathfrak p)$
by Dualizing Complexes, Lemma \ref{dualizing-lemma-torsion-change-rings}.
Thus we conclude by Algebra, Lemma \ref{algebra-lemma-characterize-zero-local}.
\end{proof}
